% GENERATED FILE. Edit canonical lessons, then run npm run generate:books.
% canonical-source-hash: fnv1a64:d65bf88f58f2ad82
% canonical-lessons: BN-C501-bisheshogyo, BN-C501-upodeshta, BN-C501-onubadok, BN-C501-sthopoti, BN-C501-nirnoy-kora

\chapter{Experts}
\label{ch:bn-experts}
\hlchaptermodality{501}

\begin{chapteropening}
\textbf{By the end of this chapter:} \emph{I can name some experts and say how they help.}

Complete the last lesson of chapter 501: I can name some experts and say how they help.
\end{chapteropening}

\section[\texorpdfstring{\bn{বিশেষজ্ঞ} (bisheshôgyô)}{বিশেষজ্ঞ (bisheshôgyô)}]{\bn{বিশেষজ্ঞ} (bisheshôgyô) — an expert; a specialist}

\label{lesson:BN-C501-bisheshogyo}



\begin{quote}
Before the new one: say the Bengali for to get done, then the Bengali for to supervise.
\end{quote}

\subsection*{You'll want to know: \bn{বিশেষজ্ঞ}}
\textbf{\bn{বিশেষজ্ঞ}} — \emph{bisheshôgyô} — \textquotedblleft{}an expert; a specialist\textquotedblright{}.

\begin{quote}
\bn{আমাদের} \bn{পাড়ার} \bn{ডাক্তার} \bn{শিশুদের} \bn{বিশেষজ্ঞ।}
\end{quote}

\begin{quote}
\emph{Āmāder pāṛār ḍāktār shishuder bisheshôgyô.}
\end{quote}

\begin{quote}
The doctor in our neighbourhood is a specialist in children.
\end{quote}

\begin{cousinweb}
\textbf{\bn{বিশেষজ্ঞ}} (\emph{bisheshôgyô}) is a tatsama, Sanskrit \emph{viśeṣajña}, \textquotedblleft{}one who knows the distinctions\textquotedblright{} (Monier-Williams).
\end{cousinweb}

\begin{grammarlens}[title={What you've built}]
One more word for describing what you lived and felt.
\end{grammarlens}

\subsection*{Your turn}
\begin{itemize}
\raggedright
  \item \emph{Say it:} \emph{bisheshôgyô}
  \item \emph{Say it:} \emph{bisheshôgyô}, once more
  \item \emph{Say it:} \emph{Āmāder pāṛār ḍāktār shishuder bisheshôgyô.} — The doctor in our neighbourhood is a specialist in children.
  \item \emph{Recall:} say the Bengali for to get done, then the Bengali for to supervise, then say \emph{bisheshôgyô} again
\end{itemize}

\begin{tcolorbox}[breakable,colback=teal!4,colframe=teal!35!black,title={Before you move on}]
What does \bn{বিশেষজ্ঞ} mean? (\textquotedblleft{}An expert; a specialist\textquotedblright{}.) Say it once more.
\end{tcolorbox}
\section[\texorpdfstring{\bn{উপদেষ্টা} (upôdeshṭā)}{উপদেষ্টা (upôdeshṭā)}]{\bn{উপদেষ্টা} (upôdeshṭā) — an adviser}

\label{lesson:BN-C501-upodeshta}



\begin{quote}
Before the new one: say the Bengali for to supervise, then the Bengali for a specialist.
\end{quote}

\subsection*{You'll want to know: \bn{উপদেষ্টা}}
\textbf{\bn{উপদেষ্টা}} — \emph{upôdeshṭā} — \textquotedblleft{}an adviser\textquotedblright{}.

\begin{quote}
\bn{নতুন} \bn{ব্যবসার} \bn{জন্য} \bn{আমি} \bn{একজন} \bn{উপদেষ্টার} \bn{পরামর্শ} \bn{নিলাম।}
\end{quote}

\begin{quote}
\emph{Nôtun bæbsār jônyô āmi ekjôn upôdeshṭār pôrāmôrshô nilām.}
\end{quote}

\begin{quote}
I took an adviser's advice for the new business.
\end{quote}

\begin{cousinweb}
\textbf{\bn{উপদেষ্টা}} (\emph{upôdeshṭā}) is a tatsama, Sanskrit \emph{upadeṣṭṛ}, \textquotedblleft{}one who teaches, an adviser\textquotedblright{} (Monier-Williams).
\end{cousinweb}

\begin{grammarlens}[title={What you've built}]
One more word for describing what you lived and felt.
\end{grammarlens}

\subsection*{Your turn}
\begin{itemize}
\raggedright
  \item \emph{Say it:} \emph{upôdeshṭā}
  \item \emph{Say it:} \emph{upôdeshṭā}, once more
  \item \emph{Say it:} \emph{Nôtun bæbsār jônyô āmi ekjôn upôdeshṭār pôrāmôrshô nilām.} — I took an adviser's advice for the new business.
  \item \emph{Recall:} say the Bengali for to supervise, then the Bengali for a specialist, then say \emph{upôdeshṭā} again
\end{itemize}

\begin{tcolorbox}[breakable,colback=teal!4,colframe=teal!35!black,title={Before you move on}]
What does \bn{উপদেষ্টা} mean? (\textquotedblleft{}An adviser\textquotedblright{}.) Say it once more.
\end{tcolorbox}
\section[\texorpdfstring{\bn{অনুবাদক} (ônubādôk)}{অনুবাদক (ônubādôk)}]{\bn{অনুবাদক} (ônubādôk) — a translator}

\label{lesson:BN-C501-onubadok}



\begin{quote}
Before the new one: say the Bengali for a specialist, then the Bengali for an adviser.
\end{quote}

\subsection*{You'll want to know: \bn{অনুবাদক}}
\textbf{\bn{অনুবাদক}} — \emph{ônubādôk} — \textquotedblleft{}a translator\textquotedblright{}.

\begin{quote}
\bn{এই} \bn{বইয়ের} \bn{অনুবাদক} \bn{একজন} \bn{বিখ্যাত} \bn{কবি।}
\end{quote}

\begin{quote}
\emph{Ei bôiyer ônubādôk ekjôn bikhyātô kôbi.}
\end{quote}

\begin{quote}
The translator of this book is a famous poet.
\end{quote}

\begin{cousinweb}
\textbf{\bn{অনুবাদক}} (\emph{ônubādôk}) is Sanskrit \emph{anuvādaka}, \textquotedblleft{}repeating with explanation\textquotedblright{} (Monier-Williams); \textquotedblleft{}translator\textquotedblright{} is the modern sense.
\end{cousinweb}

\begin{grammarlens}[title={What you've built}]
One more word for describing what you lived and felt.
\end{grammarlens}

\subsection*{Your turn}
\begin{itemize}
\raggedright
  \item \emph{Say it:} \emph{ônubādôk}
  \item \emph{Say it:} \emph{ônubādôk}, once more
  \item \emph{Say it:} \emph{Ei bôiyer ônubādôk ekjôn bikhyātô kôbi.} — The translator of this book is a famous poet.
  \item \emph{Recall:} say the Bengali for a specialist, then the Bengali for an adviser, then say \emph{ônubādôk} again
\end{itemize}

\begin{tcolorbox}[breakable,colback=teal!4,colframe=teal!35!black,title={Before you move on}]
What does \bn{অনুবাদক} mean? (\textquotedblleft{}A translator\textquotedblright{}.) Say it once more.
\end{tcolorbox}
\section[\texorpdfstring{\bn{স্থপতি} (sthôpôti)}{স্থপতি (sthôpôti)}]{\bn{স্থপতি} (sthôpôti) — an architect}

\label{lesson:BN-C501-sthopoti}



\begin{quote}
Before the new one: say the Bengali for an adviser, then the Bengali for a translator.
\end{quote}

\subsection*{You'll want to know: \bn{স্থপতি}}
\textbf{\bn{স্থপতি}} — \emph{sthôpôti} — \textquotedblleft{}an architect\textquotedblright{}.

\begin{quote}
\bn{এই} \bn{সেতুর} \bn{স্থপতি} \bn{একজন} \bn{মহিলা।}
\end{quote}

\begin{quote}
\emph{Ei setur sthôpôti ekjôn môhilā.}
\end{quote}

\begin{quote}
The architect of this bridge is a woman.
\end{quote}

\begin{cousinweb}
\textbf{\bn{স্থপতি}} (\emph{sthôpôti}) is a tatsama, Sanskrit \emph{sthapati}, which Monier-Williams gives both as \textquotedblleft{}a chief\textquotedblright{} and as \textquotedblleft{}an architect, a master builder\textquotedblright{}.
\end{cousinweb}

\begin{grammarlens}[title={What you've built}]
One more word for describing what you lived and felt.
\end{grammarlens}

\subsection*{Your turn}
\begin{itemize}
\raggedright
  \item \emph{Say it:} \emph{sthôpôti}
  \item \emph{Say it:} \emph{sthôpôti}, once more
  \item \emph{Say it:} \emph{Ei setur sthôpôti ekjôn môhilā.} — The architect of this bridge is a woman.
  \item \emph{Recall:} say the Bengali for an adviser, then the Bengali for a translator, then say \emph{sthôpôti} again
\end{itemize}

\begin{tcolorbox}[breakable,colback=teal!4,colframe=teal!35!black,title={Before you move on}]
What does \bn{স্থপতি} mean? (\textquotedblleft{}An architect\textquotedblright{}.) Say it once more.
\end{tcolorbox}
\section[\texorpdfstring{\bn{নির্ণয়} \bn{করা} (nirnôy kôrā)}{নির্ণয় করা (nirnôy kôrā)}]{\bn{নির্ণয়} \bn{করা} (nirnôy kôrā) — to determine; to diagnose}

\label{lesson:BN-C501-nirnoy-kora}



\begin{quote}
Before the new one: say the Bengali for a translator, then the Bengali for an architect.
\end{quote}

\subsection*{You'll want to know: \bn{নির্ণয়} \bn{করা}}
\textbf{\bn{নির্ণয়} \bn{করা}} — \emph{nirnôy kôrā} — \textquotedblleft{}to determine; to diagnose\textquotedblright{}.

\begin{quote}
\bn{ডাক্তার} \bn{আমার} \bn{রোগ} \bn{নির্ণয়} \bn{করলেন।}
\end{quote}

\begin{quote}
\emph{Ḍāktār āmār rog nirnôy kôrlen.}
\end{quote}

\begin{quote}
The doctor diagnosed my illness.
\end{quote}

Past (he/she …): \bn{তিনি} \bn{সমস্যাটা} \bn{নির্ণয়} \bn{করলেন}, \emph{tini sômôsyāṭā nirnôy kôrlen}.

\begin{cousinweb}
\textbf{\bn{নির্ণয়} \bn{করা}} (\emph{nirnôy kôrā}) joins \bn{নির্ণয়}, Sanskrit \emph{nirṇaya}, \textquotedblleft{}complete ascertainment, decision\textquotedblright{} (Monier-Williams), to \bn{করা}, \textquotedblleft{}to do\textquotedblright{}.
\end{cousinweb}

\begin{grammarlens}[title={What you've built}]
One more word for describing what you lived and felt.
\end{grammarlens}

\subsection*{Your turn}
\begin{itemize}
\raggedright
  \item \emph{Say it:} \emph{nirnôy kôrā}
  \item \emph{Say it:} \emph{nirnôy kôrā}, once more
  \item \emph{Say it:} \emph{Ḍāktār āmār rog nirnôy kôrlen.} — The doctor diagnosed my illness.
  \item \emph{Recall:} say the Bengali for a translator, then the Bengali for an architect, then say \emph{nirnôy kôrā} again
\end{itemize}

\begin{tcolorbox}[breakable,colback=teal!4,colframe=teal!35!black,title={Before you move on}]
What does \bn{নির্ণয়} \bn{করা} mean? (\textquotedblleft{}To determine; to diagnose\textquotedblright{}.) Say it once more.
\end{tcolorbox}
